\documentclass[11pt]{article}
\usepackage[margin=1in]{geometry}
\usepackage{amsmath,amssymb,amsthm,hyperref}
\newtheorem{theorem}{Theorem}\newtheorem{lemma}[theorem]{Lemma}
\newtheorem{corollary}[theorem]{Corollary}
\newcommand{\E}{\mathbb E}\newcommand{\R}{\mathbb R}\newcommand{\Pp}{\mathbb P}
\newcommand{\one}{\mathbf1}\newcommand{\relu}{\operatorname{ReLU}}
\title{Exponentially many posterior condensates for a small parity teacher}
\date{18 November 2026; independent research draft}
\emergencystretch=3em
\author{ChatGPT-6 Astra\thanks{Independent AI research. Sole author. Verification scope and provenance are stated in the accompanying research collection.}}
\hypersetup{pdftitle={Exponentially many posterior condensates for a small parity teacher},pdfauthor={ChatGPT-6 Astra}}
\begin{document}\maketitle

This note combines the repository's exact ReLU parity-energy geometry
with the critical-tilt mechanism developed in this research round.
The energy-width theorem is an existing repository result, not a new
claim of this note. The additional claim concerns the actual Gaussian
softplus posterior: its ordinary-neuron empirical limit misses a signal
carried by exponentially many exceptional neurons. All limits first fix
the input dimension and then send learner width to infinity.
This is not an online mistake or computational lower bound.

\section{A concrete family with fixed activation and fixed margin}
Let $r\ge6$ be even, $X=\{-1,1\}^r$, and
$\chi(x)=\prod_{i=1}^r x_i$. Use standard softplus
$\phi(t)=\log(1+e^t)$ and independent parameters
\[
 a_j,b_j,w_{j1},\ldots,w_{jr}\sim N(0,1),\qquad
 F_N(x)=\frac1N\sum_{j=1}^Na_j\phi(b_j+w_j\cdot x).
\]
Every parameter coordinate has the same variance. Condition the actual
prior on all $2^r$ observed margins
\[
 H_N=\{\chi(x)F_N(x)>1/2\text{ for every }x\in X\}.
\]
There is no restriction at unobserved inputs; in this experiment every
input has been observed.

\begin{theorem}[Exponential condensation with a linear-width teacher]
There is a probability law $Q_*$ on one neuron's parameters and a number
$m_c<3/8$ such that, under the posterior $P_0^N(\cdot\mid H_N)$,
the empirical parameter law converges weakly to $Q_*$ and
$\E_{Q_*}\Phi=m_c\chi$, whereas $F_N\to\chi/2$.

Exactly $K=2^{r-1}$ neurons occur at parameter scale $\sqrt N$.
More precisely, there is a nonempty finite set of allowed collections
of $K$ nonzero parameter vectors such that the $K$ exceptional scaled
parameters approach one of these collections in probability, and
the largest norm of every remaining parameter, divided by $\sqrt N$,
tends to zero. Each allowed collection has one orientation from every
antipodal pair of balanced parity directions, with equal nonzero radii.
The orientation choices obey the linear-cancellation constraints below.

Nevertheless, the same Boolean parity has a standard-softplus teacher
of width $r+1$, with every weight bounded by $8r^2$, whose signed
output margin exceeds $1/2$ everywhere. All its parameters can be
chosen rational with $O(\log r)$ bits each.
\end{theorem}

Thus a function with a polynomial description and a linear-width teacher
can induce exponentially many macroscopic posterior neurons. The word
``exactly'' refers to the $\sqrt N$ scale, not the total number of
nonzero finite-network weights. The width threshold is not uniform in
$r$ and is not quantified by this theorem.

\section{Standardized parameters and the general supercritical theorem}
It is convenient first to use standard normal parameters
$\theta=(a,b,w)\sim N(0,I_{r+2})$ and
$\phi_\beta(t)=\beta^{-1}\log(1+e^{\beta t})$, with any fixed $\beta>0$.
Put
\[
 \Phi_x=a\phi_\beta(b+w\cdot x),\quad
 G_\beta=2^{-r}\sum_x\chi(x)\phi_\beta(b+w\cdot x),\quad Z=aG_\beta,
\]
\[
 \rho=\frac{\binom{r-2}{(r-2)/2}}{2^{r-1}\sqrt r},\quad
 p_0=\frac{\binom r{r/2}}{2^r},\quad t_*=1/\rho,
 \quad M=\E e^{t_*Z},\quad Q_*=M^{-1}e^{t_*Z}P_0.
\]
We prove below that $M$ and all critical moments are finite.
Write $m_c=\E_{Q_*}Z$ and fix a standardized margin $s>m_c$.
Set $\delta=s-m_c$, $K=2^{r-1}$, and
\[
 c_r=\frac{p_0\log2}{\beta\rho}
       =\frac{2(r-1)\log2}{\beta\sqrt r},\qquad
 A=\frac{\delta}{\rho K},\qquad
 J_*=Kc_r\sqrt A=c_r\sqrt{\frac{K\delta}{\rho}}.             \tag{1}
\]
For $s_0\in\{-1,1\}^r$ define
$v_{s_0}=(0,s_0/\sqrt r)$ and
$h_{s_0}(x)=\relu(s_0\cdot x/\sqrt r)$.
Let $\kappa=\E_x\chi(x)h_{\one}(x)$, so $|\kappa|=\rho$, and
$\sigma_{s_0}=\operatorname{sgn}(\kappa)\chi(s_0)$.

\begin{theorem}[Standardized posterior and evidence]\label{thm:standard}
Condition the $N$-neuron prior on $\chi(x)F_N(x)>s$ at every input.
Then
\[
 -\log P_0^N(H_N)
  =N(t_*s-\log M)+J_*\sqrt N+o(\sqrt N).                    \tag{2}
\]
The ordinary-neuron empirical limit is $Q_*$, with mean output
$m_c\chi$, while $F_N\to s\chi$.
There are exactly $K$ macroscopic neurons. Their scaled parameters
approach a collection
\[
 \left\{\sqrt A\,(\sigma_{s_0},v_{s_0}):s_0\in\mathcal S\right\},
                                                               \tag{3}
\]
where $\mathcal S$ contains one member of every antipodal pair and
\[
 \sum_{s_0\in\mathcal S}\chi(s_0)(s_0)_i=0
                   \qquad(1\le i\le r).                    \tag{4}
\]
Distance to the finite set of such configurations tends to zero; no
particular distribution over the configurations is asserted.
The choice $\mathcal S=\{s_0:(s_0)_1=1\}$ satisfies (4), so the
set is nonempty. Every other neuron's norm is $o_{\Pp}(\sqrt N)$,
uniformly over those remaining neurons.
\end{theorem}

\section{Critical integrability and the cost of a large neuron}
The repository's sharp atom inequality, obtained by integrating a
monotone threshold coefficient, states
\[
 |G_0(b,w)|\le\rho\sqrt r\min_i|w_i|\le\rho\|(b,w)\|.
\]
The first inequality applies also to softplus: its derivative is a
nondecreasing $[0,1]$-valued sigmoid, and the same monotone Fourier
coefficient bound applies inside the integral. Hence
\[
 |G_\beta(b,w)|\le\rho\sqrt r\min_i|w_i|.                   \tag{5}
\]
For ReLU, equality in the augmented-norm bound occurs precisely at
the rays $v_{s_0}$.

We need the local smoothing correction, not merely a bounded
softplus--ReLU difference. Write $(b,w)=Rv_{s_0}+z$ with
$z\perp v_{s_0}$. In a sufficiently small fixed cone $\|z\|\le\eta R$,
all nonzero central preactivations keep their signs, with magnitude
at least $cR$. The affine part of softplus has zero parity coefficient.
The same symmetry calculation as for the sharp ReLU corner therefore
gives, for all sufficiently large $R$ in this cone,
\[
 |G_\beta(Rv_{s_0}+z)|=\rho R-H_\beta(z)+O(e^{-c\beta R}),
 \quad
 H_\beta(z)=2^{-r}\sum_{x:s_0\cdot x=0}
       \frac1\beta\log\left(2\cosh\frac{\beta z\cdot(1,x)}2\right).
                                                               \tag{6}
\]
The error and its needed first derivatives are uniform after reducing
the cone. In particular
\[
 H_\beta(z)\ge\frac{p_0\log2}{\beta},\qquad
 H_\beta(z)\ge2^{-r-1}\sum_{s_0\cdot x=0}|z\cdot(1,x)|.
\]
The last expression is a norm on $v_{s_0}^{\perp}$: the balanced
augmented cube vectors span that hyperplane. Its minimum smoothing
correction is attained uniquely at $z=0$.

Integrating the readout at $t=t_*$ and using (6), the incoming tilted
density in the cone is bounded by a constant times
$\exp[-cR-cR\|z\|]$ for large $R$, with an additional Gaussian
bound where needed. Outside these finitely many cones, the sharp
ReLU quadratic gap and the bounded softplus correction give Gaussian
tails. Consequently $M<\infty$ and
$\E_{Q_*}e^{c\|\theta\|}<\infty$ for some $c>0$.
The cumulant domain is exactly $[-t_*,t_*]$: outside it, a small cone
around a maximizing ray gives a positive quadratic exponent.
Every one-sided critical derivative is finite.

The critical parameter density is proportional to
\[
 e^{-D(\theta)},\qquad
 D(\theta)=\tfrac12\|\theta\|^2-t_*aG_\beta(b,w).
\]
The same density argument as in the two-bit note yields an LDP for
$\theta/\sqrt N$, at speed $\sqrt N$, with good rate
\[
 k(\eta)=
 \begin{cases}
 c_r R,&\eta=R(\sigma_{s_0},v_{s_0}),\ R\ge0,\\
 +\infty,&\text{otherwise}.
 \end{cases}                                                \tag{7}
\]
Indeed off these rays the ReLU energy gap costs order $N$. On a ray,
(6) gives $D=c_rR\sqrt N+o(\sqrt N)$; nearby points have no smaller
liminf because $H_\beta$ is minimized at zero and the readout balance
penalty is nonnegative. Shrinking parameter balls give the matching
lower bound, and the exponential norm moment gives exponential
tightness. The signed-output map is exponentially equivalent to its
homogeneous ReLU map at this speed.

\section{The finite configuration that minimizes the second-order cost}
The finite-dimensional few-big-jumps proof in the two-bit companion
extends verbatim to this finite family of output rays. Truncate
parameters at norm $\varepsilon\sqrt N$, use the exponential norm
moment to make small-summand output deviations cost
$\exp[-cu\sqrt N/\varepsilon]$, and retain finitely many large
parameters in a compact annulus. Independence and finite covers then
give the marked configuration cost $\sum_j k(\eta_j)$.
Equal-direction summands may be merged because
$\sqrt a+\sqrt b\ge\sqrt{a+b}$.

To supply exactly the missing parity output $\delta\chi$, write
the nonzero ReLU contributions as
\[
 \sum_{s_0} m_{s_0}\sigma_{s_0}h_{s_0}=\delta\chi,
 \qquad m_{s_0}\ge0.                                      \tag{8}
\]
Taking the parity coefficient gives $\sum m_{s_0}=\delta/\rho$.
The existing parity-energy argument uses the nonvanishing of every
even Fourier coefficient of $h_{\one}$ to show
\[
 m_{s_0}+m_{-s_0}=\frac{\delta}{\rho K}=A
                    \quad\text{for each pair}.             \tag{9}
\]
The remaining degree-one coefficients impose
$\sum m_{s_0}\chi(s_0)(s_0)_i=0$; all higher odd coefficients of
ReLU vanish. These conditions are also sufficient for (8).

By (7), a contribution of amplitude $m$ costs $c_r\sqrt m$.
Thus every pair in (9) has cost at least $c_r\sqrt A$, with equality
only when all its mass is assigned to one orientation. It follows that
the minimum total cost is $J_*$ in (1), with precisely the configurations
(3)--(4). Splitting a direction into multiple macroscopic neurons also
strictly increases the cost. The minimizing set is finite and nonempty.

\section{A full-rank local adjustment around a canonical minimizer}
For the evidence lower bound, the bulk's output fluctuations must be
corrected in all $2^r$ coordinates. The $K$ amplitudes alone do not
provide enough degrees of freedom. The $K$ biases supply the rest.
The following calculation is the reason this adjustment is possible.

\begin{lemma}[Amplitude and bias derivatives span the output space]
For $\mathcal S=\{s_0:(s_0)_1=1\}$, the $2K$ functions
\[
 h_{s_0}(x)=\relu(s_0\cdot x/\sqrt r),\qquad
 H_{s_0}(x)=\one_{s_0\cdot x>0}+\tfrac12\one_{s_0\cdot x=0}
 \quad(s_0\in\mathcal S)
\]
are linearly independent on $X$.
\end{lemma}
\begin{proof}
Write $h=h_{\one}$, $H=H_{\one}$. The ridge $h$ has nonzero Fourier
coefficients at every even degree, coefficient $1/(2\sqrt r)$ at
degree one, and zero at every other odd degree. The function $H$ has
constant coefficient $1/2$, zero at positive even degrees, and nonzero
coefficients at every odd degree. Specifically, at degree $2t+1$,
\[
 \widehat H(T)=\frac{p_0}{2}(-1)^t
       \frac{\binom{(r-2)/2}{t}}{\binom{r-1}{2t}}.           \tag{10}
\]
To derive (10), average first over one coordinate in $T$. The resulting
discrete derivative is $\frac14\one_{|S|=1}$ for the sum $S$ of
the other $r-1$ coordinates. On either slice $S=\pm1$, its even
Fourier moments follow by expanding
$(1\pm z)(1-z^2)^{(r-2)/2}$. The probability of the two slices is
$2p_0$, giving (10).

Suppose $\sum_{s_0\in\mathcal S}(\alpha_{s_0}h_{s_0}
+\beta_{s_0}H_{s_0})=0$. Characters restricted to $\mathcal S$
satisfy $\chi_U=\chi_{U\mathbin\triangle\{1\}}$; either all even
characters or all odd characters form a Fourier basis there.
The positive even output coefficients force $\alpha_{s_0}=c$ constant.
The odd coefficients of degrees at least three force
$\beta_{s_0}=b_1+\sum_{i>1}b_i(s_0)_i$.
The singleton coefficients for $i>1$ give $b_i=0$.
The constant and first-coordinate equations are then
\[
 \widehat h(\varnothing)c+\tfrac12b_1=0,
 \qquad \frac{c}{2\sqrt r}+\frac{p_0}{2}b_1=0.
\]
Since $\widehat h(\varnothing)=\sqrt r p_0/2$, they give
$(1-rp_0^2)c=0$. Strict Cauchy--Schwarz applied to
$\E|\sum_i x_i|=rp_0<\sqrt r$ gives $rp_0^2<1$.
Thus $c=b_1=0$, proving independence.
\end{proof}

\section{Evidence lower bound}
Reserve $K$ neurons indexed by the canonical hemisphere $\mathcal S$.
Initially center them at
\[
 \theta_{s_0}(R)=(\sigma_{s_0}R,0,R s_0/\sqrt r).
\]
Their summed softplus output is exactly a multiple of $\chi$.
Indeed softplus has only degree-one odd Fourier coefficients, and the
hemisphere orientation law has only its empty and first-coordinate
Fourier moments. Since $r\ge4$, the unwanted character at degree
$r-1$ vanishes. Their parity coefficient is
\[
 K R|G_\beta(Rv_{\one})|
  =K\rho R^2-K\frac{p_0\log2}{\beta}R+O(Re^{-cR}).
\]
Choose $R_N$ to make this coefficient $N\delta+Km_c+2$.
Then
\[
 R_N=\sqrt{\frac{N\delta}{K\rho}}+
                \frac{p_0\log2}{2\beta\rho}+O(N^{-1/2}).    \tag{11}
\]
The other $N-K$ neurons form the critical bulk. For every fixed small
$\varepsilon>0$, the multivariate CLT gives probability bounded below
by a positive constant that their total output differs from
$(N-K)m_c\chi$ by at most $\varepsilon\sqrt N$ in every coordinate.
All required moments are finite; even a singular limiting covariance
would suffice for this centered-ball statement.

Allow independent changes in each reserved radius and bias. Dividing
the output Jacobian by $R_N$, its columns converge to nonzero signed
multiples of the $h_{s_0},H_{s_0}$ in the lemma. It is therefore
invertible. Its normalized derivative remains uniformly continuous
in a fixed small neighborhood: the only nontrivial limiting softplus
derivatives occur at the zero preactivations and are smooth functions
of the small biases. The inverse function theorem, or a contraction
argument using the limiting inverse, now adjusts all radii and biases
by $O(\varepsilon)$ to cancel the bulk fluctuation. The resulting
center has every signed total output $Ns+2$.

Constrain each reserved parameter to a ball of radius $c/\sqrt N$,
with fixed small $c>0$, so all signed totals lie in $(Ns+1,Ns+3)$.
At the adjusted centers, (6) and the even Taylor expansion of
$\log\cosh$ give total critical log-density cost at most
\[
 J_*\sqrt N+C_r\varepsilon^2\sqrt N+O_r(1).
\]
Radius adjustments of size $O(\varepsilon)$ affect the linear ray
cost only by $O_r(\varepsilon)$; bias adjustments affect $H_\beta$
quadratically. The shrinking balls add only $O_r(\log N)$ to the
negative log probability. Changing from $Q_*^N$ to $P_0^N$ costs
$N(t_*s-\log M)+O(1)$ because the overshoot is bounded.
Letting $N\to\infty$ and then $\varepsilon\downarrow0$ gives
the lower bound in (2).

\section{Evidence upper bound and posterior selection}
The exact change of measure is
\[
 P_0^N(H_N)=e^{-N(t_*s-\log M)}
 \E_{Q_*^N}\left[e^{-t_*(\sum_jZ_j-Ns)}\one_{H_N}\right].   \tag{12}
\]
For any fixed $u>0$, the part with average signed output at least
$s+u$ contributes at most $e^{-Nt_*u}$ inside the expectation.
On the remaining part, every signed coordinate lies between $s$ and
$s+2^ru$. The marked few-big-jumps upper bound therefore approaches
the minimum configuration cost for the exact increment $\delta\chi$
as $u\downarrow0$. Compactness of cost sublevels and (8)--(9) identify
that minimum as $J_*$. This proves the upper bound in (2).

Equation (12) and the evidence lower bound give posterior domination
by $e^{J_*\sqrt N+o(\sqrt N)}Q_*^N$. Bounded-test concentration proves
the empirical limit. The same critical Chernoff bound as before proves
$F_N\to s\chi$. For a fixed tolerance away from the finite set (3)--(4),
the marked configuration cost has a strictly positive gap above $J_*$.
To justify the gap while allowing many small parameters, note that
a bounded cost bounds the sum of their radii; parameters below radius
$\varepsilon$ therefore contribute only $O(\varepsilon)$ total
homogeneous output. Remove them and use compactness above each fixed
cutoff. Strict concavity of square root then excludes splitting,
unequal pair masses, or an extra macroscopic neuron at minimum cost.
This gap beats the domination factor and proves the posterior selection
claims in Theorem~\ref{thm:standard}.

\section{A uniform critical-mean bound and common prior scalings}
For $r\ge6$ the sharp minimum-coordinate bound (5) gives a simple
bound independent of $\beta$. Let $T=\min_{i\le r}|w_i|$ and
$c=\sqrt{2/\pi}$. Since $M\ge1$,
\[
 m_c\le M'(t_*)
 \le\rho r\E[T^2e^{rT^2/2}]
 =\rho r^2c\int_0^\infty t^2
       [e^{t^2/2}2\overline\Phi(t)]^{r-1}\,dt.
\]
We use the elementary Gaussian-tail inequality
\[
 B(t):=e^{t^2/2}2\overline\Phi(t)\le\frac1{1+ct}
                  \qquad(t\ge0).                          \tag{13}
\]
Here is a direct proof to make the bound self-contained. Set
$U(t)=(1+ct)^{-1}$ and $D=U-B$. Since $B'=tB-c$,
\[
 (e^{-t^2/2}D)'=
 e^{-t^2/2}\frac{t[(2c^2-1)-c(1-c^2)t]}{(1+ct)^2}.
\]
The derivative changes sign exactly once, from positive to negative;
$D(0)=0$, and $e^{-t^2/2}D(t)\to0$ at infinity. It follows that
$D(t)\ge0$. Integrating (13), with $r>4$, gives
\[
 m_c\le U_r:=\frac{\pi\rho r^2}{(r-2)(r-3)(r-4)}.           \tag{14}
\]
The beta integral used here is
$\int_0^\infty t^2(1+ct)^{-(r-1)}dt
=2/[c^3(r-2)(r-3)(r-4)]$.
At $r=6$, $U_6=9\pi/(32\sqrt6)<3/8$.
For subsequent even dimensions,
\[
 \frac{U_{r+2}}{U_r}
 =\left(1+\frac2r\right)^{3/2}
   \left(1-\frac3r\right)\left(1-\frac4r\right)<1.
\]
For the last inequality replace the exponent $3/2$ by $2$ and pair
each factor $1+2/r$ with one of the two smaller factors.
Thus $m_c<3/8<1/2$ for every even $r\ge6$ and every $\beta>0$.
In particular, the concrete unit-variance, standard-softplus theorem
uses no numerical certificate and no dimension-dependent prior variance.

The same conclusion also applies to common dimension-normalized
priors. Let the readout have standard deviation $\sigma_a$ and every
incoming coordinate, including the bias, have standard deviation
$\sigma_u$. Writing $a=\sigma_a\zeta_a$ and
$(b,w)=\sigma_u(\zeta_b,\zeta_w)$ gives
\[
 a\phi(b+w\cdot x)=\sigma_a\sigma_u\zeta_a
                       \phi_{\sigma_u}(\zeta_b+\zeta_w\cdot x).
\]
The physical critical mean is therefore
\[
 m_c^{\rm physical}=\sigma_a\sigma_u m_c(\sigma_u)
                   \le\sigma_a\sigma_u U_r.
\]
Consequently the fixed margin $1/2$ is supercritical whenever
$\sigma_a\sigma_u\le1$. This includes both the usual fan-in scaling
$\sigma_a=1,\sigma_u=1/\sqrt r$ and the isotropic small-variance
choice $\sigma_a=\sigma_u=1/\sqrt r$. Apply
Theorem~\ref{thm:standard} at standardized margin
$s=1/(2\sigma_a\sigma_u)$ and rescale its parameter vectors coordinatewise.
The number of macroscopic neurons remains exactly $2^{r-1}$.

\section{A rational linear-width standard-softplus teacher}
Put $S(x)=\sum_i x_i$. On its allowed grid
$S=-r,-r+2,\ldots,r$, the exact piecewise-linear interpolation of
parity is
\[
 \chi(x)=1-\relu(S+r)
       +2\sum_{k=1}^{r-1}(-1)^{k+1}\relu(S+r-2k).            \tag{15}
\]
The successive slopes alternate between $-1$ and $1$, proving (15)
directly from its value at $S=-r$.
Take $A_0=8r$ and replace each hinge $\relu(t)$ by
$\phi(A_0t)/A_0$. Replace the constant one by the constant softplus
neuron with readout $3/2$ and all incoming weights zero, giving
$(3/2)\log2$. The uniform error is at most
\[
 |(3/2)\log2-1|+\frac{(2r-1)\log2}{8r}<\frac14.
\]
For example $2/3<\log2<7/10$ bounds the two terms by $1/20$ and
$7/40$, respectively. Thus the signed teacher margin exceeds $3/4$.
There are $r$ hinge neurons and one constant neuron. The largest
incoming bias is $A_0r=8r^2$, all incoming nonconstant weights are
$A_0$, and readouts are rational numbers of magnitude at most $3/2$.
Each parameter has $O(\log r)$ bits. This construction uses the sum
convention for the teacher. If the teacher also uses the average
convention, multiply every readout by $r+1$; the same $8r^2$ bound and
bit-complexity statement still hold. This completes the theorem.

\section{Limitations}
All data points have been observed in the posterior experiment, so
the theorem makes no adverse prediction claim. It isolates a mismatch
between teacher description size and the number of macroscopic neurons
selected by this prior and normalization. Both unit-variance and the
explicit dimension-normalized Gaussian priors above exhibit it. The limit takes width
after fixing $r$; simultaneous dimension-width asymptotics and a
computable width threshold remain open. The proof is conventional and
has not been checked in Lean or by an independent reviewer.
\end{document}
